\documentclass[12pt]{article}
\usepackage[utf8]{inputenc}
\usepackage[T1]{fontenc}
\usepackage[margin=1in]{geometry}
\usepackage{setspace}
\usepackage{xcolor}
\usepackage[colorlinks=true, urlcolor=blue!60!black, linkcolor=blue!60!black]{hyperref}
\setlength{\parindent}{2em}
\setlength{\parskip}{1.25em}
\renewcommand{\baselinestretch}{1.0}
\begin{document}
\begin{center}
{\LARGE \textit{West Coast Hotel Co. v. Parrish}, 300 U.S. 379 (1937)\par}
\vspace{0.8em}
{\large Case Brief\par}
\end{center}
\vspace{1em}

\section{Memory Jogger}

In a landmark 5--4 decision that ended the \textit{Lochner} era of economic substantive due process, the Supreme Court overruled \textit{Adkins v. Children's Hospital} and upheld Washington's minimum-wage law for women as a reasonable exercise of state police power under the Fourteenth Amendment's Due Process Clause (\href{https://www.courtlistener.com/opinion/102784/west-coast-hotel-co-v-parrish/}{CourtListener}).

\section{Detailed Case Facts}

In 1913, the Washington State Legislature enacted a ``Minimum Wages for Women'' statute (Laws of 1913, ch. 174) declaring that the welfare of the State demanded protecting women and minors from labor conditions having a pernicious effect on their health and morals (\href{https://caselaw.findlaw.com/court/us-supreme-court/300/379.html}{FindLaw}; \href{https://en.wikipedia.org/wiki/West_Coast_Hotel_Co._v._Parrish}{Wikipedia}). The statute made it unlawful to employ women at wages inadequate for their maintenance and authorized an administrative body---originally the Industrial Welfare Commission, and later the Industrial Welfare Committee (comprising state labor and industrial officials, including the Supervisor of Women in Industry)---to investigate workplace conditions, convene conferences of employers, employees, and public representatives, and issue binding minimum-wage orders (\href{https://caselaw.findlaw.com/court/us-supreme-court/300/379.html}{FindLaw}; \href{https://fedsoc.org/case/west-coast-hotel-company-v-parrish}{Federalist Society}). The statute also allowed special licenses permitting lower wages for workers incapable of full service and granted underpaid employees a civil cause of action to recover the full statutory minimum wage notwithstanding any contract to work for less (\href{https://www.studicata.com/case-briefs/case/west-coast-hotel-co-v-parrish}{Studicata}). Acting under this authority, the state committee fixed a minimum wage of \$14.50 for a 48-hour workweek for hotel chambermaids (\href{https://fedsoc.org/case/west-coast-hotel-company-v-parrish}{Federalist Society}; \href{https://en.wikipedia.org/wiki/West_Coast_Hotel_Co._v._Parrish}{Wikipedia}).

Elsie Parrish worked as a chambermaid at the Cascadian Hotel in Wenatchee, Washington, owned and operated by the West Coast Hotel Company, from August 1933 until May 1935 (\href{https://supreme.justia.com/cases/federal/us/300/379/}{Justia}; \href{https://en.wikipedia.org/wiki/West_Coast_Hotel_Co._v._Parrish}{Wikipedia}). Throughout her employment, the hotel paid her less than the state-mandated \$14.50 weekly minimum (\href{https://supreme.justia.com/cases/federal/us/300/379/}{Justia}). After being discharged, Parrish, joined by her husband, sued the West Coast Hotel Company to recover \$216.19 (roughly \$216)---the difference between the wages actually paid or tendered under her employment contract and the statutory minimum wage (\href{https://supreme.justia.com/cases/federal/us/300/379/}{Justia}; \href{https://en.wikipedia.org/wiki/West_Coast_Hotel_Co._v._Parrish}{Wikipedia}). The hotel did not dispute the wage arithmetic; instead, it defended on the ground that Washington's minimum-wage statute violated the Fourteenth Amendment's Due Process Clause by abridging freedom of contract, relying on the Supreme Court's decision in \textit{Adkins v. Children's Hospital}, 261 U.S. 525 (1923) (\href{https://fedsoc.org/case/west-coast-hotel-company-v-parrish}{Federalist Society}).

\section{Procedural History}

On October 17, 1935, Judge W. O. Parr of the Chelan County Superior Court entered judgment for the West Coast Hotel Company, holding Washington's 1913 minimum-wage statute unconstitutional under the controlling authority of \textit{Adkins} (\href{https://en.wikipedia.org/wiki/West_Coast_Hotel_Co._v._Parrish}{Wikipedia}; \href{https://fedsoc.org/case/west-coast-hotel-company-v-parrish}{Federalist Society}). On direct appeal, the Supreme Court of Washington unanimously reversed and directed judgment for Parrish, upholding the statute as a valid exercise of the state police power and distinguishing \textit{Adkins}, 185 Wash. 581, 55 P.2d 1083 (1936) (\href{https://supreme.justia.com/cases/federal/us/300/379/}{Justia}).

The West Coast Hotel Company appealed to the U.S. Supreme Court. Only months earlier, in \textit{Morehead v. New York ex rel. Tipaldo}, 298 U.S. 587 (1936), a 5--4 Court (with Justice Owen Roberts joining the four conservative Justices) had invalidated New York's minimum-wage law for women, stating that the petitioner had not asked the Court to reconsider or overrule \textit{Adkins} (\href{https://caselaw.findlaw.com/court/us-supreme-court/300/379.html}{FindLaw}). In \textit{West Coast Hotel}, by contrast, the question whether \textit{Adkins} should be reconsidered was squarely presented (\href{https://caselaw.findlaw.com/court/us-supreme-court/300/379.html}{FindLaw}). The Supreme Court heard oral argument on December 16--17, 1936, and issued its decision affirming the Washington Supreme Court on March 29, 1937 (\href{https://supreme.justia.com/cases/federal/us/300/379/}{Justia}).

\section{Judicial Votes}

The Supreme Court affirmed the judgment of the Supreme Court of Washington by a \textbf{5--4} vote (\href{https://www.courtlistener.com/opinion/102784/west-coast-hotel-co-v-parrish/}{CourtListener}; \href{https://fedsoc.org/case/west-coast-hotel-company-v-parrish}{Federalist Society}):
\begin{itemize}
\item \textbf{Majority (5):} Chief Justice Charles Evans Hughes (author), joined by Associate Justices Louis D. Brandeis, Harlan Fiske Stone, Benjamin N. Cardozo, and Owen J. Roberts.
\item \textbf{Dissent (4):} Associate Justice George Sutherland (author), joined by Associate Justices Willis Van Devanter, James Clark McReynolds, and Pierce Butler (the bloc colloquially known as the ``Four Horsemen'').
\end{itemize}
Notably, Justice Roberts cast his conference vote to uphold Washington's statute on December 19, 1936---leaving the Court temporarily deadlocked 4--4 while Justice Stone was absent due to illness---nearly two months before President Franklin D. Roosevelt announced his Judicial Procedures Reform Bill (``court-packing plan'') in February 1937 (\href{https://en.wikipedia.org/wiki/West_Coast_Hotel_Co._v._Parrish}{Wikipedia}; \href{https://www.stevevladeck.com/p/42-the-switch-in-time-that-saved}{Steve Vladeck}). Once Stone returned and joined the majority, the March 29, 1937 announcement became popularly known as ``the switch in time that saved nine'' (\href{https://www.stevevladeck.com/p/42-the-switch-in-time-that-saved}{Steve Vladeck}).

\section{Holding}

\textbf{Judgment:} Affirmed. Washington's minimum-wage statute for women does not violate the Due Process Clause of the Fourteenth Amendment, and \textit{Adkins v. Children's Hospital}, 261 U.S. 525 (1923), is expressly overruled (\href{https://supreme.justia.com/cases/federal/us/300/379/}{Justia}; \href{https://www.studicata.com/case-briefs/case/west-coast-hotel-co-v-parrish}{Studicata}).

\textbf{Rule Announced:} The Constitution does not guarantee an absolute ``freedom of contract''; rather, liberty under the Fourteenth Amendment's Due Process Clause is subject to reasonable legislative regulation adopted in the interests of the community to protect health, safety, morals, and public welfare (\href{https://constitutioncenter.org/the-constitution/supreme-court-case-library/west-coast-hotel-co-v-parrish}{National Constitution Center}). So long as an economic or labor regulation has a reasonable relation to a proper legislative purpose and is neither arbitrary nor capricious, it satisfies due process, and a legislature may address evils piecemeal by targeting classes of workers---such as women---where the need is deemed clearest (\href{https://caselaw.findlaw.com/court/us-supreme-court/300/379.html}{FindLaw}).

\textbf{Holding vs. Dicta:} The \emph{holding} is that (1) a state may constitutionally enact minimum-wage legislation for women to protect their health and prevent exploitative wages, (2) \textit{Adkins} is overruled, and (3) omitting men from the statute does not render it an arbitrary discrimination under the Fourteenth Amendment (\href{https://supreme.justia.com/cases/federal/us/300/379/}{Justia}). Chief Justice Hughes's broader reflections---taking judicial notice of Great Depression relief rolls to characterize sub-living wages as a taxpayer ``subsidy for unconscionable employers'' and suggesting that legislative judgment controls whenever the wisdom of economic policy is debatable---operate as \emph{reasoning and expansive dicta} that laid the doctrinal foundation for near-total judicial deference to all economic regulations (including gender-neutral wage laws and non-labor price controls), later cemented in \textit{United States v. Carolene Products Co.}, 304 U.S. 144 (1938), and \textit{Williamson v. Lee Optical}, 348 U.S. 483 (1955) (\href{https://constitutioncenter.org/the-constitution/supreme-court-case-library/west-coast-hotel-co-v-parrish}{National Constitution Center}; \href{https://caselaw.findlaw.com/court/us-supreme-court/300/379.html}{FindLaw}).

\section{Analysis of Opinions}

\subsection{Majority Opinion (Chief Justice Hughes)}

Chief Justice Hughes began by explaining why \textit{Adkins} was properly before the Court despite \textit{Morehead}: whereas New York had not asked the Court to overrule \textit{Adkins} in \textit{Morehead}, the Washington Supreme Court had sustained a statute closely paralleling the D.C. law invalidated in \textit{Adkins}, compelling direct reconsideration (\href{https://caselaw.findlaw.com/court/us-supreme-court/300/379.html}{FindLaw}).

Turning to the merits, Hughes dismantled the constitutional premise of \textit{Lochner}-era substantive due process. ``The Constitution does not speak of freedom of contract,'' he wrote; ``[i]t speaks of liberty and prohibits the deprivation of liberty without due process of law'' (\href{https://constitutioncenter.org/the-constitution/supreme-court-case-library/west-coast-hotel-co-v-parrish}{National Constitution Center}). Constitutional liberty is ``liberty in a social organization,'' not an ``absolute and uncontrollable liberty,'' and regulation that is ``reasonable in relation to its subject and is adopted in the interests of the community is due process'' (\href{https://constitutioncenter.org/the-constitution/supreme-court-case-library/west-coast-hotel-co-v-parrish}{National Constitution Center}). Surveying decades of precedent---including \textit{Muller v. Oregon}, 208 U.S. 412 (1908), \textit{O'Gorman \& Young v. Hartford Fire Insurance Co.}, 282 U.S. 251 (1931), and \textit{Nebbia v. New York}, 291 U.S. 502 (1934)---Hughes emphasized that legislatures enjoy broad discretion to regulate employer-employee relations and private contracts, and that ``every possible presumption is in favor of'' a statute's validity (\href{https://caselaw.findlaw.com/court/us-supreme-court/300/379.html}{FindLaw}).

Hughes offered two reinforcing justifications for Washington's law. First, echoing \textit{Muller}, he stressed the State's special interest in protecting women's health against oppressive working conditions and unequal bargaining power (\href{https://supreme.justia.com/cases/federal/us/300/379/}{Justia}). Second, and more innovatively, he rejected the \textit{Lochner}-era baseline that treated market wages as natural and state intervention as redistributive: taking judicial notice of Great Depression relief demands, Hughes observed that when employers pay less than a living wage, ``[w]hat these workers lose in wages the taxpayers are called upon to pay,'' and ``[t]he community is not bound to provide what is in effect a subsidy for unconscionable employers'' (\href{https://constitutioncenter.org/the-constitution/supreme-court-case-library/west-coast-hotel-co-v-parrish}{National Constitution Center}). Finally, Hughes dismissed the hotel's equal-protection-style objection that the statute covered only women, holding that a legislature ``is free to recognize degrees of harm'' and need not regulate every class simultaneously (\href{https://caselaw.findlaw.com/court/us-supreme-court/300/379.html}{FindLaw}). Accordingly, the Court expressly overruled \textit{Adkins} (\href{https://www.studicata.com/case-briefs/case/west-coast-hotel-co-v-parrish}{Studicata}).

\subsection{Dissenting Opinion (Justice Sutherland)}

Justice Sutherland, joined by Justices Van Devanter, McReynolds, and Butler, delivered a forceful dissent grounded in constitutional fixity, judicial duty, and formal legal equality (\href{https://caselaw.findlaw.com/court/us-supreme-court/300/379.html}{FindLaw}). Pointedly rebuking the majority's reliance on Great Depression conditions---and implicitly chastising Justice Roberts's shift---Sutherland declared that ``the meaning of the Constitution does not change with the ebb and flow of economic events'' (\href{https://constitutioncenter.org/the-constitution/supreme-court-case-library/west-coast-hotel-co-v-parrish}{National Constitution Center}). While constitutional text applies to new factual circumstances, ``[t]he judicial function is that of interpretation; it does not include the power of amendment under the guise of interpretation'' (\href{https://caselaw.findlaw.com/court/us-supreme-court/300/379.html}{FindLaw}). Treating judicial self-restraint as a license to abdicate constitutional enforcement, Sutherland warned, converts binding constitutional mandates into ``mere moral reflections'' (\href{https://caselaw.findlaw.com/court/us-supreme-court/300/379.html}{FindLaw}).

On the substance of liberty of contract, Sutherland reaffirmed \textit{Adkins}: Washington's statute was ``essentially arbitrary'' because it forced employers to pay wages based not on the value of the services rendered, hours worked, or workplace hazards, but on the extraneous personal subsistence needs of the worker---ignoring the ``just equivalence'' between pay and work (\href{https://caselaw.findlaw.com/court/us-supreme-court/300/379.html}{FindLaw}). Moreover, Sutherland objected to the statute's paternalistic sex classification: following the Nineteenth Amendment and the demise of common-law coverture, ``[w]omen today stand upon a legal and political equality with men'' and should neither be relegated to an inferior contractual class nor effectively barred from competing with men for jobs paying lower wages (\href{https://constitutioncenter.org/the-constitution/supreme-court-case-library/west-coast-hotel-co-v-parrish}{National Constitution Center}).

\section{Examples: Future Applications}

\begin{enumerate}
\item \textbf{Same-Side (Upheld --- State Minimum Wage for All Adult Workers):} A state enacts a gender-neutral minimum wage of \$15 per hour for all adult retail and hospitality workers, and a restaurant chain challenges the law under the Fourteenth Amendment's Due Process Clause as a deprivation of freedom of contract. \emph{Why:} Under \textit{West Coast Hotel}, freedom of contract is not absolute, and a wage floor reasonably related to worker welfare and preventing taxpayer-funded public assistance for underpaid workers easily survives deferential due process review.

\item \textbf{Same-Side (Upheld --- State Price Floor on Agricultural Commodities):} A state legislature imposes a minimum retail price on dairy products to stabilize farm incomes and prevent destructive price wars, and a discount grocer sues alleging infringement of economic liberty. \emph{Why:} Building directly on \textit{Nebbia} and \textit{West Coast Hotel}, courts presume economic regulations constitutional and will not second-guess the legislature's policy judgment so long as the measure is not arbitrary or capricious.

\item \textbf{Opposite-Side (Invalidated --- Economic Regulation Targeting a Fundamental Right):} A state passes a statute prohibiting private employers from hiring any person who has publicly criticized the state governor, defending the law as a general regulation of employment contracts under \textit{West Coast Hotel}. \emph{Why:} \textit{West Coast Hotel} curtailed substantive due process scrutiny of ordinary \emph{economic} regulations; it did not authorize states to abridge enumerated constitutional guarantees such as First Amendment political speech, which trigger strict scrutiny rather than rational-basis deference.

\item \textbf{Opposite-Side (Invalidated --- Modern Sex-Specific Protective Wage Law):} A state today reenacts a minimum-wage law applying exclusively to female workers, citing \textit{West Coast Hotel}'s reliance on \textit{Muller v. Oregon} and piecemeal legislative discretion. \emph{Why:} Although \textit{West Coast Hotel} rejected a due process challenge to a women-only wage law in 1937, modern Equal Protection jurisprudence (developed in \textit{Craig v. Boren} and \textit{United States v. Virginia}) subjects facial sex classifications to intermediate scrutiny and forbids paternalistic generalizations about women's workplace capacities---making Justice Sutherland's equality critique prescient under contemporary doctrine.

\item \textbf{Fence-Sitter (Borderline --- Purely Protectionist Occupational Licensing):} A state requires independent African-style hair braiders---who use no chemicals or heat---to complete 1,500 hours of unrelated cosmetology school tuition at a cost of \$20,000, and braiders challenge the law as an irrational deprivation of economic liberty designed solely to shield licensed salons from competition. \emph{Why:} Under \textit{West Coast Hotel}'s highly deferential standard, courts routinely uphold licensing laws if any conceivable health or safety rationale exists; yet several modern federal courts of appeals have split over whether naked economic protectionism divorced from any plausible public-welfare purpose crosses the line into the ``arbitrary or capricious'' regulation that even \textit{West Coast Hotel} forbids.
\end{enumerate}

\section{Critique}

\textbf{Originalist and Libertarian Perspectives:} Classical-liberal and libertarian constitutional scholars---such as Randy Barnett, Richard Epstein, and David Bernstein---criticize \textit{West Coast Hotel} for inaugurating a near-total judicial abdication over economic liberties (\href{https://www.heritage.org/the-constitution/report/economic-liberty-and-the-constitution-introduction}{Heritage Foundation}; \href{https://digitalcommons.wcl.american.edu/cgi/viewcontent.cgi?article=1061&context=aulr}{American University Law Review}). From this vantage point, the decision replaced meaningful policing of arbitrary ``class legislation'' with an unprincipled double standard that treats property and contract rights as second-class liberties while permitting rent-seeking interest groups to capture state legislatures (\href{https://www.heritage.org/the-constitution/report/economic-liberty-and-the-constitution-introduction}{Heritage Foundation}; \href{https://insight.dickinsonlaw.psu.edu/cgi/viewcontent.cgi?article=1185&context=dlr}{Dickinson Law Review}). Critics also note the irony that Hughes's opinion relied heavily on \textit{Muller}-era paternalism regarding women, whereas Sutherland's dissent defended women's formal legal equality after the Nineteenth Amendment (\href{https://www.heritage.org/the-constitution/report/economic-liberty-and-the-constitution-introduction}{Heritage Foundation}).

\textbf{Progressive, Historical, and Doctrinal Perspectives:} Progressive and doctrinal scholars celebrate \textit{West Coast Hotel} as a vital correction to the anti-democratic judicial activism of the \textit{Lochner} era, crediting Hughes with exposing the fallacy of treating common-law market baselines as constitutionally neutral (\href{https://www.nyulawreview.org/wp-content/uploads/2018/08/NYULawReview-76-5-Friedman.pdf}{NYU Law Review}). By recognizing that sub-living wages shift subsistence costs onto taxpayers, Hughes reframed labor regulation as a legitimate public-welfare choice rather than impermissible wealth redistribution (\href{https://constitutioncenter.org/the-constitution/supreme-court-case-library/west-coast-hotel-co-v-parrish}{National Constitution Center}). Meanwhile, revisionist legal historians---most prominently Barry Cushman---have challenged the popular ``switch in time'' narrative, demonstrating from docket chronology and doctrinal evolution (especially \textit{Nebbia}) that Justice Roberts's vote reflected internal doctrinal development rather than political capitulation to Roosevelt's court-packing threat (\href{https://chicagounbound.uchicago.edu/cgi/viewcontent.cgi?article=5058&context=uclrev}{University of Chicago Law Review}; \href{https://www.stevevladeck.com/p/42-the-switch-in-time-that-saved}{Steve Vladeck}).

\textbf{Underlying Competing Values:} The clash in \textit{West Coast Hotel} embodies enduring tensions in constitutional law: democratic self-governance versus countermajoritarian judicial review; individual contractual autonomy versus collective social welfare; formal legal neutrality versus substantive economic equality; and constitutional adaptability to changing economic realities versus fixed constitutional meaning.

\section{Key Quotations}

\begin{quote}
``What is this freedom? The Constitution does not speak of freedom of contract. It speaks of liberty and prohibits the deprivation of liberty without due process of law. In prohibiting that deprivation, the Constitution does not recognize an absolute and uncontrollable liberty. Liberty in each of its phases has its history and connotation. But the liberty safeguarded is liberty in a social organization which requires the protection of law against the evils which menace the health, safety, morals, and welfare of the people. Liberty under the Constitution is thus necessarily subject to the restraints of due process, and regulation which is reasonable in relation to its subject and is adopted in the interests of the community is due process.''
\end{quote}
--- Chief Justice Charles Evans Hughes, Majority Opinion, 300 U.S. at 391 (\href{https://constitutioncenter.org/the-constitution/supreme-court-case-library/west-coast-hotel-co-v-parrish}{National Constitution Center}; \href{https://www.courtlistener.com/opinion/102784/west-coast-hotel-co-v-parrish/}{CourtListener}).

\begin{quote}
``The exploitation of a class of workers who are in an unequal position with respect to bargaining power and are thus relatively defenseless against the denial of a living wage is not only detrimental to their health and well being, but casts a direct burden for their support upon the community. What these workers lose in wages the taxpayers are called upon to pay. The bare cost of living must be met. \ldots\ The community is not bound to provide what is in effect a subsidy for unconscionable employers.''
\end{quote}
--- Chief Justice Charles Evans Hughes, Majority Opinion, 300 U.S. at 399 (\href{https://constitutioncenter.org/the-constitution/supreme-court-case-library/west-coast-hotel-co-v-parrish}{National Constitution Center}; \href{https://www.courtlistener.com/opinion/102784/west-coast-hotel-co-v-parrish/}{CourtListener}).

\begin{quote}
``It is urged that the question involved should now receive fresh consideration, among other reasons, because of `the economic conditions which have supervened'; but the meaning of the Constitution does not change with the ebb and flow of economic events.''
\end{quote}
--- Justice George Sutherland, Dissenting Opinion, 300 U.S. at 402 (\href{https://constitutioncenter.org/the-constitution/supreme-court-case-library/west-coast-hotel-co-v-parrish}{National Constitution Center}; \href{https://www.courtlistener.com/opinion/102784/west-coast-hotel-co-v-parrish/}{CourtListener}).

\begin{quote}
``The judicial function is that of interpretation; it does not include the power of amendment under the guise of interpretation. To miss the point of difference between the two is to miss all that the phrase `supreme law of the land' stands for and to convert what was intended as inescapable and enduring mandates into mere moral reflections.''
\end{quote}
--- Justice George Sutherland, Dissenting Opinion, 300 U.S. at 404 (\href{https://caselaw.findlaw.com/court/us-supreme-court/300/379.html}{FindLaw}; \href{https://www.courtlistener.com/opinion/102784/west-coast-hotel-co-v-parrish/}{CourtListener}).

\begin{quote}
``Women today stand upon a legal and political equality with men. There is no longer any reason why they should be put in different classes in respect of their legal right to make contracts; nor should they be denied, in effect, the right to compete with men for work paying lower wages which men may be willing to accept.''
\end{quote}
--- Justice George Sutherland, Dissenting Opinion, 300 U.S. at 411--412 (\href{https://constitutioncenter.org/the-constitution/supreme-court-case-library/west-coast-hotel-co-v-parrish}{National Constitution Center}; \href{https://www.courtlistener.com/opinion/102784/west-coast-hotel-co-v-parrish/}{CourtListener}).

\end{document}
